\documentclass[../main.tex]{subfiles}
\begin{document}
%==========================================TIÊU ĐỀ TẬP========================================
\begin{table}[h]
  \begin{adjustwidth}{-1cm}{}
    \begin{tabular}{>{\centering\arraybackslash}p{6cm}|>{\raggedright\arraybackslash}p{12.5cm}}
      \multirow{5}{6cm}
      % Cột bên trái: ÉPISODE và số tập
      {\\[-40pt]
        \flaregothic\fontsize{20pt}{20pt}\selectfont \centering
        \color{eptype}ÉPISODE\color{black}\\
        \burbank\fontsize{72pt}{72pt}\selectfont
        \color{epnum}111\color{black}\\[6pt]
        \cafeteria\fontsize{20pt}{20pt}\selectfont\color{epnum}(9-3)\color{black}
        \\[18pt]
        % \flaregothic\fontsize{14pt}{14pt}\selectfont
        % \color{epattr}BIENVENUE, \uline{\color{Banpire} Mel} !
        \color{black}
      }
       &
      % Dòng thứ nhất: tên tập tiếng Nhật
      {\fotskip\fontsize{18pt}{18pt}\selectfont お\ruby{悩}{なや}み\ruby{解}{かい}\ruby{決}{けつ}、\ruby{失}{しっ}\ruby{敗}{ぱい}！
      } \\[6pt]
      % Dòng thứ hai: tên tập tiếng Anh
       & {\cafeteria\fontsize{18pt}{18pt}\selectfont Advice Unwise}                                             \\[4pt]
      % Dòng thứ ba: tên tập tiếng Việt
       & {\vnmsans\fontsize{18pt}{18pt}\selectfont Nhà tư vấn "có cũng như không" Pekora}                       \\[8pt]
      % Dòng thứ tư: Nhân vật xuất hiện
       & {\caviar{\fontsize{14pt}{14pt}\selectfont Nhân vật: \emp{kuon1} \emp{ryuuhou1} \emp{pekora} \emp{rushia} \emp{towa} \emp{lamy}}}\\
            % Dòng thứ năm (nếu có): Nhân vật xuất hiện trong Omake
            & {\abv Truyện phụ:} \emp{kuon2} \emp{achan} \emp{korone}
            \\[8pt]
      % Dòng cuối cùng: Thời gian diễn ra của tập (thường sẽ là ngày phát hành)
       & {\continuum\fontsize{16pt}{16pt}\selectfont Ngày phát hành: 20/06/2021}                                \\[18pt]
    \end{tabular}
  \end{adjustwidth}
\end{table}
%=========================================TRUYỆN CHÍNH========================================
\color{black}

\txtbx[2]
{{\color{vol} \uline{Tóm tắt:}} Pekora mở ``Trung tâm Tư vấn'' nhưng ca đầu tiên của Rushia về vấn đề ngực khiến cô đuổi Pekora bằng dao, buộc nàng Thỏ phải kích hoạt bẫy hố. Ca thứ hai Towa muốn trông ác ma hơn, Pekora xúi cô dùng búa cao su đập Koron, để rồi Towa bị nàng Khuyển trói đem đi không rõ số phận. Đến lượt Lamy từ tủ lạnh chui ra và buộc Pekora phải tự tư vấn cho chính mình về thói quen nói ``peko''. Kuon và Ryuuhou - với kinh nghiệm dẫn dắt nhóm - giúp Pekora chấp nhận bản thân, rồi lần lượt tháo gỡ những trăn trở của Rushia và Towa. Sáng hôm sau, A-chan gọi cho Kuon báo một tin khiến cậu đứng hình: có một thành viên của \hll\ sắp tốt nghiệp. Một trang mới đang mở ra, và một trang cũ đang dần khép lại\dots}

Ngày nắng đầu tháng Sáu, ánh sáng mặt trời vừa nhú lên trên đỉnh nhà xuyên qua ô cửa sổ, rọi thẳng xuống phòng khách nơi tất cả mọi người đang tụ tập. Không khí có vẻ êm đềm và yên bình, cho đến khi\dots

``Trung tâm Tư vấn Pekora chính thức khai trương! Nếu bạn đang có vấn đề cần được giải đáp, thì hãy đến với chúng tôi, bạn đã tìm đúng chỗ rồi đấy!''

Nàng Thỏ đứng chễm chệ ở giữa căn phòng, hai tay giương cao lên trời, gương mặt rạng rỡ đầy tự tin. Trước mặt cô là một chiếc bàn học nhỏ xinh, giọng nói tràn đầy năng lượng vang vọng khắp từng góc phòng.

\img{9-3a}

Ngồi đối diện với cô, Kuon chỉ biết khoanh tay lại, khẽ nhíu mày nhìn cảnh tượng khó tin đang diễn ra trước mắt. Bên cạnh cậu, Rushia nằm bẹp dí trên ghế sô-pha, ánh mắt cũng chẳng giấu được sự hoài nghi nặng nề. Ở một góc khác của căn phòng, Ryuuhou ngồi xếp bằng trên chiếc ghế đẩu, tay cầm ly trà đá mát lạnh, mắt nhìn Pekora với nét thích thú pha lẫn chút nghi ngờ.

``Sao tự dưng bà ấy lại nảy ra ý định mở trung tâm tư vấn vậy nhỉ?'' nàng Pháp sư lên tiếng, gương mặt tỏ rõ sự bối rối, đôi mắt híp lại nhìn cô bạn cùng Thế hệ y như đang xem một vở hài kịch.

Nàng Trưởng nhóm nhấp nhẹ một ngụm trà, rồi cất giọng đầy tinh nghịch: ``Em thấy Pekora-san trông cứ như đang đóng phim vậy. Mà trung tâm tư vấn của chị ấy có giấy phép kinh doanh không đấy?''

Cậu Tổng thở dài não nề, lắc đầu bất lực: ``Mấy hôm trước anh thấy em ấy lướt mạng đọc tin, thấy trên báo có nhiều người đang tìm chỗ tư vấn tâm lý, thế là hùa theo làm theo người ta thôi.''

Nàng Pháp sư chớp mắt ngạc nhiên, rồi nhún vai, giọng nói vẫn đầy ngờ vực: ``\dots Thà Kuon-senpai làm tư vấn viên thì em thấy còn hợp lý hơn á-nanodesu. Bình thường bà ấy toàn gây chuyện, tự nhiên lại muốn đi tư vấn cho người khác là sao không biết?''

Ryuuhou gật gù như đồng tình hoàn toàn, thêm vào: ``Em cũng nghĩ giống Rushia-san. Onii trông trưởng thành và đàng hoàng hơn Pekora-san nhiều. Chưa kể Pekora-san còn chưa kịp hỏi câu nào thì đã bắt đầu lúng túng rồi kìa.''

``Anh cũng chẳng biết phải làm sao bây giờ\dots''

Pekora đang háo hức chuẩn bị đón khách hàng đầu tiên của mình. Cô đột nhiên quay đầu sang phía Kuon, giọng nói tràn đầy hứng khởi: ``Anh! Anh có chuyện gì muốn hỏi không, peko!?''

Cậu vẫn giữ được thái độ bình tĩnh, đáp lại một cách từ tốn nhưng vẫn không giấu được sự nghi ngờ, giọng cậu vẫn thản nhiên: ``Hiện tại thì anh chưa có chuyện gì cần nhờ giúp đỡ cả.''

Ryuuhou cười khúc khích tiếp lời: ``Em cũng không có điều gì phải hỏi cả.''

Nhận được câu trả lời chẳng mấy vui vẻ từ hai anh em, nàng Thỏ lập tức chuyển sang đối tượng còn lại: ``Vậy thì\dots\ bà kia! Có gì thì cứ nói thẳng ra đi!''

Rushia nhìn Pekora một lúc, ánh mắt vẫn giữ nguyên vẻ chán chường, rồi lạnh lùng thốt lên, giọng nói đều đều nhưng sát thương không kém: ``Muốn biết làm thế nào để ngực to hơn.''

\img{9-3b}

Khoảnh khắc đó, cả căn phòng chìm vào im lặng tuyệt đối. Ngay cả tiếng ve sầu kêu ngoài cửa sổ cũng như bị bối cảnh làm cho ngừng lại.

Pekora chớp mắt liên tục, như không tin vào những gì vừa lọt vào tai, miệng há hốc ra không khép lại được. Kuon cũng sững người trong giây lát, bàn tay đang cầm cốc nước đột nhiên khựng lại giữa không trung. Ryuuhou thì suýt nữa phun hết ngụm trà trong miệng ra, vội đưa tay bịt miệng lại, mắt tròn thao nhìn Rushia vừa bất ngờ vừa thán phục.

``Mới ca đầu tiên mà sao lại đặt ra câu hỏi khó thế!'' cả ba người cùng đồng thanh thốt lên trong sự kinh ngạc. Cô nàng mắt hai màu lắc đầu, giọng nói bất lực nhưng vẫn nở được nụ cười: ``Rushia-san ơi, chị hỏi câu đó thì tới em cũng có thể hỏi được kia kìa. Chọn câu nào không hỏi mà sao hỏi khó thế!''

Nhưng dù sao đi nữa, với tư cách là một tư vấn viên, Pekora đã tuyên bố sẽ giúp đỡ khách hàng của mình, nên cô đành phải cố gắng giải quyết.

``Được rồi, để anh suy nghĩ xem\dots'' cô đưa tay lên cằm trầm ngâm, trán cô nhíu lại y như một vị thám tử đang điều tra một vụ án phức tạp, rồi nghiêm chỉnh hỏi, giọng cố gắng giữ vững sự tự tin: ``Đã thử uống bổ sung khoáng chất chưa?''

Rushia thở dài, giọng nói như thể đã nghe câu này hàng ngàn lần: ``Làm đủ cách rồi mà chẳng ăn thua.''

``Vậy thử mát-xa thì sao?''

``Cũng đã thử rồi,'' cô đáp lại, đôi mắt dần dần tối sầm lại, bắt đầu mất kiên nhẫn, hai hàm răng nghiến chặt vào nhau.

Ryuuhou chứng kiến cảnh tượng đó, liền ghé sát tai Kuon mà thì thầm: ``Onii, em thấy Pekora-san sắp toang rồi. Nếu Rushia-san nổi giận lên thì chẳng dễ dỗ đâu.''

Kuon khoanh tay lại, mồ hôi vã ra như tắm: ``Xem ra ca này căng thẳng quá nhỉ\dots''

Pekora cũng bắt đầu cảm thấy có gì đó không ổn. Cô nhìn cô nàng Pháp sư một lúc, rồi hắng giọng, gật gù nói ra mấy câu đầy triết lý, giọng cố tỏ ra là người từng trải: ``À thì\dots\ mỗi người sinh ra ai cũng như vậy thôi. Cuộc đời vốn dĩ đã là thế rồi mà.''

Kuon liếc sang phía cô, hỏi bằng giọng nửa đùa nửa thật, tay chống cằm: ``Bí quá không nghĩ ra lời khuyên nào khác à?''

Ryuuhou thêm vào, giọng nói đầy sự thương hại: ``Chị nói vậy khác gì bảo Rushia-san chấp nhận số phận luôn đâu.''

Câu trả lời của nàng Thỏ dĩ nhiên chẳng làm hài lòng khách hàng. Rushia bỗng dưng nheo mắt đầy nguy hiểm, chất giọng vốn mềm mại bỗng trở nên sắc bén như lưỡi dao: ``Bà đang nói ý gì hả?''

Cả Kuon, Ryuuhou lẫn Pekora đều đứng hình, chẳng dám hó hé gì. Tình hình càng trở nên tồi tệ khi Pekora lại còn nói thêm những lời vòng vo, giọng run cầm cập: ``Thì\dots\ ý tôi là\dots\ bà cứ chấp nhận thực tế đi\dots''

Nhưng thay vì im lặng chịu trận, Pekora lại quyết định đổ thêm dầu vào lửa.

Cô hắng giọng, chắp tay ra sau lưng, nở nụ cười tự tin rồi nói với giọng đầy động viên, y như đang cổ vũ một đứa trẻ: ``Cứ cùng nhau vượt qua khó khăn thôi, peko-rin♪''

*BỘP!*

Kuon chưa kịp phản ứng gì thì Rushia đã bật dậy như một con mãnh thú, rút ngay con dao quen thuộc của mình ra từ trong áo rồi không nói không rằng mà đuổi theo Pekora, mặt đỏ bừng bừng vì giận, mắt sáng rực như hai ngọn lửa đang cháy.

``\textbf{Bà vừa nói là sẽ giải quyết được vấn đề của tôi mà! Đồ dối trá!!}'' nàng Pháp sư hét lên, giọng nói vang vọng khắp cả căn phòng.

\img{9-3c}

``\textbf{Tôi cũng chẳng còn cách nào khác để giúp bà đâu!}'' nàng Thỏ hét lên, hoảng sợ chạy thục mạng quanh phòng, vừa chạy vừa ngoái đầu nhìn cô bạn cùng Thế hệ đang đuổi sát nút phía sau.

Ryuuhou nhìn cảnh tượng hỗn loạn diễn ra, lắc đầu rồi quay sang Kuon, giọng nói hài hước nhưng vẫn giữ được nụ cười trên môi: ``Onii, trung tâm tư vấn này chẳng khác gì một buổi diễn hài kịch. Em thấy Pekora-san nên đổi tên thành `Trung tâm giải trí' thì hợp lý hơn nhiều.''

``Nơi này có bao giờ nghiêm túc đâu mà,'' ông anh cười chát đáp lại.

Rushia với gương mặt đỏ bừng vì tức giận, vẫn cầm dao đuổi theo Pekora không chịu dừng, trong khi Kuon chỉ đứng nhìn cảnh tượng loạn xạ với nụ cười chát chúa, tay xoa thái dương: ``Nói thế cũng khác gì bảo em ấy mãi mãi bị lép vế luôn.''

Bỏ ngoài tai những lời nhận xét vừa trêu chọc lại thẳng thắn của cậu Tổng, Pekora vội giật lấy sợi dây thừng treo lơ lửng trên trần nhà. Ngay lập tức, một cái bẫy hố mở ra dưới sàn nhà, và Rushia lọt thẳng xuống dưới với một tiếng hét ngắn ngủi.

\img{9-3d}

Từ sâu dưới cái hố, Rushia la hét ầm ĩ vì quá tức giận, cảm thấy bị đồng đội phản bội: ``\textbf{Đồ lừa đảo! Tôi ghét bà!!}''

Pekora thở hổn hển, đứng ở miệng hố, lấy tay quệt đi giọt mồ hôi trên trán rồi quay sang than thở với Kuon, mặt đỏ bừng vì vừa chạy vừa căng thẳng than trách: ``Ngay từ đầu bà ấy đã đưa cho em một ca khó như thế, làm sao em giúp được! Em đã cố gắng hết sức rồi mà!''

Kuon chỉ lắc đầu ngao ngán, đưa tay xoa trán: ``Nói thế cũng phũ phàng quá. Em còn non kinh nghiệm trong việc này lắm.''

Ryuuhou đứng dậy, bước đến gần miệng hố, nhìn xuống Rushia đang đứng dưới đó với vẻ thương hại nhưng không giấu được nụ cười: ``Rushia-san, chị có ổn không vậy? Em thấy Pekora-san vẫn chưa có câu trả lời nào cho chị đâu.''

Pekora chống hai tay lên hông, mặt đầy bất mãn, thở dốc nhưng vẫn cố gắng giữ vững vẻ tự tin: ``Vậy thì để xem ca tiếp theo có dễ hơn không!''

Cô vỗ tay một cái, rồi tuyên bố với giọng tràn đầy nhiệt huyết: ``Tiếp theo!''

Ryuuhou nhìn sang Kuon, giọng nói đầy sự dự đoán: ``Onii, em cá là ca tiếp theo chắc chắn lại có chuyện gì đó xảy ra. Pekora-san mà cứ thế này thì trung tâm tư vấn của Pekora-san sẽ sớm đóng cửa thôi.''

\ct

Người tiếp theo cần tư vấn là Towa, cô đến với một nỗi băn khoăn đã tồn tại từ lâu với chính mình.

Nàng Ác ma bước vào với gương mặt đầy phiền muộn, hai tay cứ xoa vặn vào nhau trước ngực, dáng vẻ cứ như sắp phải thú nhận một tội lỗi lớn lao. Cô ngồi xuống ghế đối mặt với Pekora, hít một hơi thật sâu, rồi mới bộc bạch điều khiến cô trăn trở: ``Em muốn mình trông giống một ác ma hơn.''

Điều cô nhận lại chỉ là ánh mắt bối rối và ngạc nhiên từ vị tư vấn viên Pekora. Nàng Thỏ chớp mắt liên tục, cố gắng hiểu thông tin vừa nghe phải: ``Sao lại đến hỏi một con thỏ với câu hỏi kiểu này chứ\dots'' cô lẩm bẩm, tay đưa lên cằm suy nghĩ mông lung.

Còn Kuon thì chỉ khoanh tay, nhìn Pekora với ánh mắt khó hiểu, giọng nói đầy giục giã: ``Thì cứ cố gắng giúp em ấy xem sao nào.''

Ryuuhou ngồi ở góc phòng, tay vẫn giữ ly trà như cũ, nhìn Towa với vẻ thích thú. Cô bé nhấp nhẹ một ngụm, rồi nhận xét với giọng điệu trêu đùa: ``Towa-san mà còn muốn trở thành ác ma hơn nữa sao? Em cứ nghĩ chị ấy đã đủ đáng sợ rồi đấy.''

\img{9-3e}

Nàng Thỏ nhăn mặt suy nghĩ một lúc, trán nhíu lại như một thám tử đang vắt óc tìm ra manh mối cho vụ án khó nhằn, rồi đập mạnh tay xuống bàn: ``Vậy thì em thử chơi khăm một ai đó xem sao!''

Ngay lập tức, cậu Tổng nhắc lại, giọng cậu như thể đã thấy tình huống này quá nhiều lần: ``Không phải em ấy đã làm thế nhiều lần rồi sao? Cái lần em ấy mới đến văn phòng lần đầu\footnote{Xem lại {\flaregothic ÉPISODE 65}, quyển số Năm.} với cả lần em gài Shion-chan\footnote{Xem lại {\flaregothic ÉPISODE 108}, quyển số Tám.} ấy để bị hốt đi nữa.''

Cô em gái của cậu nhướng mày, nhìn Pekora với vẻ đầy nghi ngờ: ``Pekora-san, chị thấy Towa-san làm vậy vẫn chưa đủ hả? Chị ấy mà ác ma hơn nữa thì em sợ cả công ty sẽ sụp đổ mất.''

``Thế thì phải làm nhiều hơn nữa chứ! Đó là việc cơ bản nhất mà một ác ma nào cũng hay làm!'' Pekora nhấn mạnh, tay giơ lên trời như một diễn giả đang nhiệt tình thuyết trình, mắt sáng rực như vừa khám phá ra một chân lý vĩ đại.

Towa nghe vậy lập tức phấn chấn hẳn, tung nắm đấm lên trời tỏ vẻ hoàn toàn đồng tình, gương mặt bừng sáng niềm vui: ``Đúng ha!'' cô reo lên, như thể vừa tìm ra đáp án cho câu đố khó nhằn nhất.

Pekora phấn khởi đưa cho đàn em một cây búa đồ chơi sáng loáng, cười khoái chí, giọng nói cứ như một thương gia đang chào hàng: ``Cầm lấy cái búa cao su này, khi gặp người tiếp theo thì cứ gõ mạnh vào đầu họ nha!''

\img{9-3f}

Nàng Ác ma hí hửng cầm lấy cây búa, đưa lên ngang tầm mắt để ngắm nghía như một đứa trẻ vừa nhận được quà Giáng sinh, trong khi Kuon thì nhăn mặt với vẻ lo lắng: ``Anh thấy cách này không ổn tí nào đâu\dots''

Ryuuhou nhìn cây búa trên tay Towa, lắc đầu với giọng đầy chưng hửng: ``Pekora-san, lần trước chị bắt Rushia-san phải chấp nhận số phận, lần này lại bảo Towa-san đi đánh người. Chị đang đi tư vấn hay đang xúi giục người khác dùng bạo lực vậy?''

Nhưng trước khi kịp có ai nói thêm gì, thì có một người bước vào từ bên ngoài. Kuon chỉ thoáng thấy cánh tay đeo chiếc vòng cổ màu đỏ, cậu lập tức nhận ra đó là ai. Cậu giật mình, vội đưa tay ra muốn ngăn cản, nhưng đã quá muộn.

Pekora hứng chí hét lên, giọng vang vọng khắp cả căn phòng: ``\textbf{Ngay bây giờ, peko!}''

Towa nhanh chóng giơ cây búa lên cao, còn Kuon thì hoảng hốt đưa tay ra cản, mặt cậu tái xanh: ``Ơ, khoan đã, đó là\dots''

*CÁCH!*

``\dots Koro-chan mà\dots'' cậu vừa nói xong thì cũng là lúc cây búa cao su giáng mạnh xuống đầu vị khách vừa bước vào. Một tiếng động khô khốc vang lên trong phòng, và mọi thứ lập tức chìm vào im lặng.

\img{9-3g}

Ryuuhou đang nhấp trà thì bỗng khựng lại, mắt mở toang vì kinh hoàng: ``Trời ơi, Towa-san vừa đánh Korone-san kìa! Thế là hỏng bét rồi!''

Cả Pekora và Towa đều đứng sững người, mặt xanh lét, mồ hôi lạnh túa ra như tắm. Riêng Kuon thì chỉ biết thở dài, lắc đầu với nụ cười chát chúa, giọng nói đầy bất lực: ``Kiểu này đúng là xúi dại chơi ngu có thưởng rồi còn gì.''

Korone từ từ ngoảnh đầu lại, đôi mắt vốn trong veo bỗng chốc tối sầm như một vực thẳm không đáy, cả người cô tỏa ra một luồng sát khí đáng sợ đến mức không khí xung quanh như ngưng đọng lại. Cô ngân dài một tiếng ``Ồ\~{}?'', kéo từng âm tiết một cách chậm rãi khiến sống lưng Towa lạnh toát, từng sợi tóc trên đầu cô lập tức dựng đứng lên.

Nàng Ác ma hoảng sợ lùi lại vài bước, tay run cầm cập, làm rơi cây búa xuống sàn với tiếng cạch nhẹ, miệng lắp bắp: ``E-E-Em xin lỗi chị! Em chỉ muốn mình trông ác ma hơn thôi mà!'' cô nói, giọng cứ như sắp khóc đến nơi, mắt ướt nhòe vì quá sợ hãi.

Nàng Khuyển từ từ cúi đầu, nhếch mép cười với vẻ đầy ẩn ý, đôi mắt vẫn nhìn chằm chằm vào Towa như một con thú săn mồi đang quan sát con mồi của mình. Dường như cô chẳng mấy tin vào những gì đàn em cô vừa nói, và cô cũng chẳng có ý định tha thứ cho chuyện này một cách dễ dàng.

Towa vội quay ngoắt người về phía sau, hỏi Pekora với giọng đầy cầu cứu: ``Đ-Đúng không nào, Pekora-chan!?'' nhưng khi cô quay lại nhìn nàng Thỏ thì\dots\ chẳng thấy bóng dáng Pekora đâu cả.

Ryuuhou chớp mắt một cái, rồi mới nhận ra\dots\ cô ấy đã chui tọt xuống cái hố thỏ dưới sàn nhà từ lúc nào không hay! Cô bé chỉ tay xuống cái hố, giọng vừa mỉa mai vừa bất lực: ``Onii, Pekora-san trốn mất rồi kìa. Đúng là một tư vấn viên `có tâm' có khác!''

\img{9-3h}

Kuon nhìn thấy cảnh tượng đó, không nhịn được mà bật cười lớn, tay chỉ về cái hố: ``Trốn tránh trách nhiệm luôn à?''

Towa cũng tức giận hét lên, giọng đầy phẫn nộ: ``\textbf{Cái hố thỏ này là cái gì thế đây!?}'' cô giậm mạnh chân xuống sàn, nhưng chỉ nhận lại tiếng vọng vang lên từ dưới hố.

Cậu Tổng nhún vai, nói ra một câu châm biếm: ``Có chơi phải có chịu chứ, vị tư vấn viên ơi!''

Nàng Trưởng nhóm thêm vào, giọng có vẻ hả hê nhưng vẫn không giấu được sự lo lắng: ``Towa-san, chị mà không chạy đi sớm thì Korone-san sẽ biến chị thành món đồ chơi mới của chị ấy đấy.''

Towa chưa kịp phản ứng gì thêm thì bỗng cảm thấy có thứ gì đó siết chặt quanh người mình. Cô ngạc nhiên cúi xuống: ``Ơ\dots'' và trước khi kịp hiểu chuyện gì đang xảy ra, cô đã bị Korone trói toàn thân bằng một sợi dây thừng từ lúc nào không hay, siết chặt đến mức cô chẳng thể nào cử động được.

Nàng Khuyển cúi xuống nhìn cô, vẫn giữ được nụ cười ngọt ngào nhưng đầy nguy hiểm, giọng nói vẫn mềm mại nhưng đủ khiến ai nghe cũng phải lạnh sống lưng: ``Chúng ta cũng nên ngồi lại nói chuyện với nhau một chút nhỉ, Towa-chan?''

Rồi cô quay sang Kuon, nhẹ nhàng hỏi với giọng đều đều: ``Kuon-senpai, anh có thể cho em mượn em ấy một lúc chứ?''

Kuon nhìn Towa, rồi lại nhìn sang Korone, sau đó chỉ thở dài lắc đầu, giọng cậu cứ như một lời tiễn biệt: ``Chơi ai không chơi, lại dám đùa đúng với Koro-chan thế này thì\dots\ thôi, hết cứu! Xin vĩnh biệt cụ rồi!'' cậu vẫy tay như thể đang tiễn một người bạn lên đường đi xa.

Ryuuhou đứng dậy, nhìn Towa bị kéo đi, giọng đầy thương hại nhưng vẫn giữ được nụ cười: ``Towa-san, em sẽ cầu nguyện cho chị. Lần sau đừng có nghe lời Pekora-san nữa nhé.''

Towa hoảng loạn vùng vẫy, hét lên thảm thiết, tiếng la hét của cô vang vọng khắp căn phòng: ``\textbf{Khôôôôôôông!!!}''

\img{9-3i}

Tiếng la hét của nàng Ác ma tội nghiệp vẫn vang vọng trong phòng khi Korone kéo cô đi vào bóng tối của hành lang, tiếng bước chân của nàng Khuyển đều đặn, còn Towa thì vẫn gào thét cầu cứu, nhưng tiếng dần dần nhỏ lại rồi biến mất hẳn. Một lúc sau, khi mọi thứ đã trở nên yên ắng, cái đầu của Pekora mới rụt rè ló lên từ trong cái hố, mắt đảo quanh như một con thỏ đang cảnh giác, xác nhận rằng họ đã đi xa hẳn, rồi mới thở phào nhẹ nhõm, phủi phủi chiếc váy của mình: ``Hơ\dots\ hú hồn hú vía, peko\dots''

Kuon chống cằm, nhìn cô với giọng đầy châm biếm, cứ như một người thầy đang dạy dỗ học trò: ``Lại đi tư vấn cho người khác làm điều dại dột rồi em ơi\dots''

Ryuuhou bước lại gần miệng hố, nhìn xuống Pekora với vẻ thương hại nhưng vẫn không giấu được nụ cười trên môi: ``Pekora-san, em thấy trung tâm tư vấn của chị nên đổi tên thành `Trung tâm Tạo Nghiệp' thì hợp lý hơn đấy.''

Pekora chẳng còn biết làm gì ngoài thở dài ngao ngán, tay ôm trán sau những chuyện kinh hoàng vừa xảy ra, cảm thấy đầu mình sắp nổ tung vì quá căng thẳng.

Dù sao đi nữa, cô vẫn cố gắng giữ được phong thái chuyên nghiệp của một vị tư vấn viên. Cô gõ nhẹ ngón tay lên mặt bàn, nói một cách uể oải, giọng vẫn còn run run: ``Người tiếp theo\dots''

Rồi cô nằm dài ra bàn, thở dài, chờ đợi người tiếp theo bước vào phòng.

Ryuuhou ngồi xuống ghế cạnh Kuon, nhấp ngụm trà cuối cùng, rồi nhận xét với giọng đầy dự đoán: ``Onii ơi, em có cảm giác ca tiếp theo chắc chắn lại có chuyện gì đó xảy ra rồi đấy\dots''

\ct

Mọi người cứ tưởng sẽ lại có thêm người đến tìm Pekora để nhờ tư vấn, thì bất ngờ, một bóng người đột nhiên vội vã nhảy ra từ trong tủ lạnh, người run cầm cập và thốt lên:
``Trời ơi\dots\ rét quá đi thôi\~{}!!''

\gif{9-3j}{1}{81}

Cả hai người trong phòng đều nhìn cô với ánh mắt đầy khó hiểu. Ryuuhou cũng bị giật mình, suýt nữa đánh rơi cốc trà trong tay, mắt mở to tròn nhìn Lamy với vẻ không thể tin được.

``Cái peko gì thế này?'' Pekora cụp tai xuống, ngơ ngác nhìn cô em gái đang co ro vì lạnh, người vẫn còn run rẩy như vừa mới bị tắm dưới vòi nước đá.

``Tự dưng lại chui vào trong tủ lạnh làm gì vậy?'' Kuon cũng toát mồ hôi với hành động kỳ quặc của cô nàng Yêu tinh xứ tuyết, tay không ngừng xoa trán.

Ryuuhou đặt ly trà xuống bàn, giọng vừa hài hước vừa ngạc nhiên: ``Lamy-san, chị đang tìm cách hạ nhiệt hay định làm đông lạnh luôn bản thân mình vậy?''

Nghe vậy, Lamy hớn hở đáp lại, giọng vẫn còn hơi run vì lạnh nhưng tràn đầy nhiệt huyết: ``Mùa hè này thì ở trong tủ lạnh mới là tuyệt nhất mà!''

Cậu Tổng càng nghe càng thấy khó hiểu, mặt nhăn nhó lại: ``Có ai dở hơi như em mà chui vào đó đâu\dots\ Trông em y như mấy chú chó mèo ở miền Nam quê anh mà hay chui vào tủ lạnh trốn nóng ấy.''

Ryuuhou thêm vào, giọng đầy mỉa mai: ``Lamy-san, em thấy chị nên xin vào làm nhân viên bảo quản lạnh thì hợp lý hơn đấy. Mà, em tưởng chị là yêu tinh xứ lạnh chứ? Tưởng chị chịu lạnh giỏi lắm mà.''

``Vả lại trời còn chưa nóng đến mức thế đâu peko,'' Pekora nói, tay còn vẫy vẫy trước mặt như thể đang xua đuổi bầu không khí lạnh lẽo tỏa ra từ người Lamy.

Cô nàng hắng giọng một tiếng, cố gắng lấy lại phong thái chuyên nghiệp của một vị tư vấn viên, rồi quay lại với chủ đề chính: ``Chắc là Lamy-chan cũng chẳng có nỗi lo nào ở trên đời này phải không peko?''

``Đúng vậy ạ!!'' Lamy đáp lại đầy nhiệt huyết, giọng nói vang vọng khắp cả căn phòng, khiến cả Kuon và Pekora đều toát mồ hôi hột vì lo ngại. Ryuuhou nhìn Lamy với vẻ thán phục, nhưng cũng chẳng giấu được nụ cười trên môi.

``Trả lời gì mà tự nhiên quá thế\dots'' Kuon lẩm bẩm một mình.

Thế rồi, Lamy chống nạnh, ưỡn ngực đầy tự tin như một vị sứ giả của mùa đông đang chuẩn bị tuyên bố điều gì đó quan trọng: ``Lamy em đây sẽ là người lắng nghe những chia sẻ từ chị!''

Cả hai người còn lại ngơ ngác nhìn nhau, rồi tự chỉ vào bản thân mình: ``Là bọn anh á?'' Kuon hỏi, giọng đầy ngạc nhiên.

Lamy lắc đầu, rồi chỉ thẳng vào nàng Thỏ tóc xanh: ``Không phải, chỉ có Pekora-senpai thôi. Kuon-senpai có nhiều kinh nghiệm sống hơn chị ấy, đó là điều hiển nhiên mà!'' cô nói, giọng đầy khẳng định như thể đó là một sự thật không thể chối cãi.

Ryuuhou gật gù đồng tình, giọng đầy hài hước: ``Onii mà làm tư vấn thì em thấy còn hợp hơn Pekora-san đấy. Thảo nào Lamy-san cũng biết rõ ai mới là người đáng tin cậy hơn.''

``Ờm\dots\ anh\dots\ cảm ơn?'' Kuon vẫn còn ngơ ngác, chưa hiểu mình đang được khen hay bị mỉa mai, để cho Lamy bắt đầu trò chuyện với vị tư vấn viên Pekora.

Pekora nghe cô em gái của cậu Tổng nói vậy, vừa tỏ vẻ nghi ngờ về khả năng tư vấn của Lamy, vừa không khỏi tự vấn bản thân: liệu mình có thực sự giỏi trong lĩnh vực này không? Cô cắn môi, nhìn đàn em với ánh mắt đầy phân vân.

Nhưng trước khi kịp phản bác hay nói thêm bất cứ điều gì, nàng Yêu tinh tuyết đã nhanh chóng ngồi xuống ghế một cách đầy khí thế, tay đặt lên mặt bàn, ra dáng của một tư vấn viên thực thụ, như một chuyên gia đang chuẩn bị cho buổi tư vấn đầu tiên của mình. ``Trước khi có thể giúp được người khác, chẳng phải Pekora-senpai cũng nên xem xét lại những vấn đề của chính mình sao?'' cô nói, giọng nghiêm túc như một vị giáo sư đang giảng bài trên giảng đường.

Ryuuhou cũng ngồi thẳng dậy, nhìn Pekora với ánh mắt tò mò: ``Pekora-san, em thấy Lamy-san nói đúng đấy. Trong nhóm của em, ai cũng phải tự nhìn nhận lại bản thân trước khi muốn giúp đỡ người khác. Chị có vấn đề gì muốn chia sẻ không?''

``Cái này thì anh phải đồng tình,'' Kuon cũng gật đầu đồng tình, tay chống cằm. ``Khi mình đã xác định được vấn đề và tìm ra cách giải quyết nó, thì mới có thể truyền đạt kinh nghiệm đó cho người khác được.''

Nàng Trưởng nhóm tiếp lời, giọng đầy từng trải như một người đã nhiều lần đứng trước ban nhạc của mình để giải quyết mâu thuẫn: ``Kinh nghiệm của em là, đôi khi mọi người chỉ cần một người lắng nghe trước khi họ có thể nghe được chính tiếng lòng của mình. Chị cứ nói ra đi, Pekora-san.''

``Phải thế thì việc tư vấn mới hiệu quả chứ!'' Lamy hớn hở nói, vừa giơ lên cục dấu chữ `rin' vừa thè lưỡi nháy mắt tỏ vẻ đáng yêu, như một cách để củng cố lập luận của mình, trông chẳng khác nào một đứa trẻ đang khoe thành tích học tập với bố mẹ.

``Quan trọng nhất còn là cách ăn nói nữa,'' Kuon tiếp lời, giọng điềm tĩnh. ``Em giống như ai đó có nickname `Hoa Tròn' ấy, cứ sau mỗi câu lại thêm chữ `peko'. Mọi người nghe cũng thấy chướng tai lắm.''

Ryuuhou cũng nhướng mày, nhìn Pekora với vẻ tò mò: ``Em cũng thắc mắc chuyện này lâu rồi. Pekora-san, chị có thể nói chuyện mà không cần thêm `peko' ở cuối câu không? Hay nó đã trở thành bản năng của chị mất rồi?''

Nghe được những lời từ hai người, nàng Thỏ dường như vừa vỡ lẽ ra một bài học cuộc sống mà cô chưa từng nghĩ đến. Cô chống cằm, trầm tư suy nghĩ thật lâu, mắt nhìn vào khoảng không vô định như đang tìm kiếm câu trả lời, rồi bất ngờ ngước nhìn hai người họ với ánh mắt tha thiết, như một đứa trẻ đang cầu xin sự giúp đỡ: ``Liệu\dots\ hai người có thể lắng nghe chị nói một chút được không?''

Lamy, Ryuuhou và Kuon đều mỉm cười, cùng đồng thanh với giọng đầy ấm áp: ``Đương nhiên là được!''

Ryuuhou cũng gật đầu, giọng nhẹ nhàng nhưng đầy động viên: ``Em cũng nghe đây, Pekora-san. Có gì thì cứ nói ra hết, bọn em không đi đâu cả.''

\ct

Pekora - người đã tự mình mở ra ``Trung tâm tư vấn'' giờ đây lại trở thành người được tư vấn. Sau khi im lặng một lúc để lấy lại tinh thần, nàng Thỏ bắt đầu chia sẻ về nỗi niềm bấy lâu của mình, giọng đầy trăn trở: ``Chị đã luôn nói có chữ `-peko' như vậy từ lúc chào đời. Nhưng chị sẽ phải nói như thế này đến bao giờ nữa peko!? Không lẽ cả đời này chị sẽ phải như thế này sao peko!?''

Cả Lamy, Kuon và Ryuuhou đều đơ người trước câu hỏi đó. Họ liếc nhìn nhau, không ai có câu trả lời nào thực sự thỏa đáng. Nàng Trưởng nhóm nhìn Pekora với vẻ thương hại, rồi lên tiếng, giọng trầm ngâm: ``Pekora-san, trong ban nhạc của em có một anh tay trống lúc nào cũng gõ theo một nhịp nhất định. Có người bảo anh ấy nên thay đổi, nhưng anh ấy bảo đó chính là linh hồn của mình. Có lẽ chữ `peko' cũng là một phần linh hồn của chị đấy.''

``Ờm\dots\ mà mỗi con người sinh ra ai cũng như vậy cả thôi. Cuộc đời vốn dĩ đã như thế rồi,'' Lamy cười gượng, cố gắng an ủi cô chị của mình, nhưng giọng nói có phần thiếu tự tin.

\img{9-3k}

Kuon thở dài đáp lại, giọng đầy bất lực: ``Giống hệt như những gì em nói với Rushia-chan hồi nãy ấy.''

Ryuuhou nhìn Lamy, lắc đầu: ``Lamy-san, chị nói thế thì y hệt Pekora-san lúc tư vấn cho Rushia-san vậy. Chị đang sao chép bài của chị ấy đấy à?''

Pekora lập tức đập mạnh tay xuống bàn, giận dữ hét lên, mặt đỏ bừng bừng: ``\textbf{Đừng có nhại lại lời của tôi chứ!? Cái này là ý gì thế hả!?}''

Kuon cười trừ đáp lại, giơ tay lên như thể đang đầu hàng: ``Cũng giống như Ru-chan, hay Lamy-chan, hay bất kỳ ai trong công ty này, kể cả anh nữa, ai mà chẳng có điểm riêng của mình.''

Ryuuhou cũng gật gù, thêm vào: ``Em nghĩ Onii nói đúng đấy. Ai cũng có cái gì đó riêng biệt, điều làm họ khác biệt với mọi người. Có khi chữ `peko' chính là thứ làm cho Pekora-san trở nên đặc biệt. Nếu chị mất đi nó, thì chị có còn là chính chị nữa không?''

Cuối cùng, Pekora chỉ có thể ngồi thừ ra, nhìn ba người đối diện với ánh mắt vừa bất lực vừa cam chịu. Những gì họ nói chẳng hề sai một chút nào. Cô nhận ra rằng mình không thể thay đổi một phần bản chất của mình chỉ vì người khác thấy nó khó chịu.

Khi bất cứ một cá thể nào được sinh ra, người đó sẽ luôn mang dòng máu và những tập tính mà gốc rễ của họ - tức là bố mẹ, hay tổ tiên - đã truyền lại. Đó là một quy luật luôn đúng trong thế giới tự nhiên. Dù là người hay là thú vật, thì ai cũng sẽ mang những dáng dấp di truyền từ đời trước. Không thể nào có một quy luật đơn giản như vậy mà lại bị phá vỡ được.

``Bởi vậy nên anh nghĩ là cứ sống là bản thân mình thì sẽ là tốt nhất,'' Kuon kết luận, giọng đầy chân thành. ``Dù em có nghịch phá hay hay bày trò chơi khăm, hay thậm chí cứ một hai câu lại là `peko!', nhưng miễn là em cảm thấy hài lòng với bản thân như vậy thì cũng không có ai cấm đoán được gì cả.''

Ryuuhou đứng dậy, bước đến bên bàn, giọng đầy đồng cảm: ``Em nghĩ điều quan trọng nhất là chị cảm thấy thế nào về bản thân mình. Dù có dùng bao nhiêu từ đệm đi nữa, miễn là chị tự tin với nó, thì nó sẽ trở thành một điểm mạnh của chị thôi. Trong ban nhạc của em\dots\ ờm, em có vài người bạn cũng có thói quen đặc biệt khi nói chuyện, và họ vẫn luôn tuyệt vời.''

Pekora ngồi ngẫm lại những gì mà Kuon và Ryuuhou vừa nói. Lời của hai anh em nhà Hikawa - một cậu chàng đã hơn hai mươi năm sống trên đời và có hiểu biết xã hội rất rộng cùng một cô nàng mười lăm tuổi với kinh nghiệm đầy mình của một trưởng nhóm ban nhạc - còn khiến cho cả Lamy, vị tư vấn viên tự xưng, cũng phải bất ngờ. Cô thở dài, giọng nói mang vẻ cam chịu nhưng cũng đầy nhẹ nhõm: ``Thôi được rồi\dots\ Chị sẽ tiếp tục nói `peko' như những gì đã làm từ trước đến nay vậy\dots''

Kuon gật gù, nở một nụ cười hài lòng: ``Thế mới đúng là em chứ.''

Lamy vỗ tay liên tục, giọng phấn khích như một đứa trẻ vừa được xem pháo hoa: ``Pekora-senpai đã chấp nhận được bản thân mình rồi! Đây là một bước tiến lớn đó!''

Ryuuhou cũng vỗ tay nhẹ, giọng đầy hài hước: ``Pekora-san đã vượt qua được nỗi sợ của chính mình. Giờ thì trung tâm tư vấn của chị có thể hoạt động trở lại rồi đấy. Mà chị có muốn em làm trợ lý không? Em cũng có kinh nghiệm trong việc này lắm.''

Pekora bây giờ chẳng còn hơi sức đâu để phản bác nữa, chỉ vẫy tay ra hiệu mời người tiếp theo vào. Cô nhắm mắt lại, cố gắng lấy lại tinh thần sau một buổi tư vấn đầy ngược đời: ``Người tiếp theo\dots''

Thực ra thì trong lòng, bản thân cô cũng chẳng muốn làm cái trò tư vấn này nữa. Cô biết mình vẫn chưa đủ sẵn sàng cho công việc này.

``Nhắc mới nhớ, Towa-chan và Rushia-chan không hiểu thế nào rồi nhỉ?'' Kuon chợt lên tiếng, mắt nhìn quanh căn phòng như thể vừa nhận ra điều gì đó bất thường.

Ryuuhou cũng ngừng lại, đưa mắt nhìn xung quanh: ``Ừ nhỉ, hai chị ấy đi tư vấn xong lâu rồi mà vẫn chưa thấy quay lại. Chẳng lẽ\dots''

Pekora mở mắt ra, chớp chớp một lúc rồi mới nhận ra\dots\ Hai người đó vẫn chưa hề quay lại. Căn phòng vắng tanh, chỉ còn có tiếng quạt trần quay đều đều và hơi thở của bốn người trong phòng.

Cả ba cùng nhìn nhau, rồi cùng nhìn về phía cánh cửa, nơi mà những người từng đi tư vấn đều đã biến mất từ lâu.

``Họ\dots\ vẫn ổn chứ nhỉ peko?'' Pekora hỏi, giọng đầy lo lắng, mắt mở to nhìn về phía cánh cửa đang đóng im lìm.

%==========================================TRUYỆN PHỤ=========================================
\omk

\begin{center}
  \uline{15 phút sau.}
\end{center}

Cả căn phòng vẫn còn trong không khí tĩnh lặng. Sau khi vừa bộc bạch xong nỗi lòng của mình, Pekora ngả lưng nhẹ về sau, khoanh tay trước ngực và hướng mắt lên trần nhà, như đang miệt mài suy tư về những vấn đề sâu sắc của cuộc đời. Cô thở một hơi dài, mí mắt cụp xuống dần.

``Liệu có ai sẽ thực sự hiểu được những gì mình làm không, peko?'' nàng Thỏ đã cảm thấy nhẹ nhõm hơn một chút, nhưng sâu trong lòng vẫn còn lảng vảng nỗi bất an và sự tự ti, như một vết nứt nhỏ nhoi trên lớp vỏ tự tin mà cô đang gắng sức vun đắp.

Ryuuhou vẫn ngồi yên ở vị trí cũ, chống cằm nhìn Pekora với giọng điệu ôn tồn và chân thành: ``Pekora-san, em nghĩ là có mà. Chỉ cần chị cứ giữ được bản thân và luôn thành thật với mọi việc mình làm, chắc chắn sẽ có người thấu hiểu. Cũng như trong ban nhạc của em vậy, có những lúc em cũng nghĩ mình chẳng có ai hiểu, nhưng cuối cùng mọi người vẫn luôn ở bên cạnh.''

Lamy nghiêng đầu, tỏ vẻ đang miệt mài suy ngẫm, rồi nở ra một nụ cười tươi: ``Nếu Pekora-san sửa được những thói quen chưa tốt của mình, mọi người chắc chắn sẽ hiểu mà! Em bảo đảm luôn!''

Kuon cũng gật gù đồng tình, bổ sung thêm với giọng trầm ổn của một người đã từng trải qua nhiều chuyện: ``Vấn đề cốt lõi không phải là người khác có hiểu hay không, mà là cách mình nhìn nhận và đối mặt với nó mới là điều quan trọng.''

Ryuuhou cũng gật đầu đồng tình, giọng nói như một người trưởng nhóm đã nhiều lần đối mặt với khó khăn trong việc dẫn dắt tập thể: ``Onii nói đúng. Đôi khi em cũng thấy các thành viên trong nhóm\dots\ ừm, những người bạn của em cũng có lúc cảm thấy bị mọi người hiểu lầm. Nhưng khi họ đủ dũng cảm để đứng lên giải thích rõ ràng, mọi người cũng sẽ sẵn sàng lắng nghe và thấu hiểu.''

Pekora im lặng. Có vẻ cô đang thực sự suy nghĩ kỹ về những lời khuyên vừa nhận được. Từ trước đến nay, cô luôn cố gắng làm hết sức mình theo cách riêng, nhưng đôi khi chính những thói quen và cách hành xử độc đáo ấy lại khiến người khác hiểu lầm. Đôi tai thỏ của cô cụp xuống, tràn đầy sự trầm tư.

Mọi thứ vẫn đang diễn ra êm đẹp trong buổi tư vấn ấy\dots

\dots cho đến khi những tiếng ồn ào từ bên ngoài hành lang vọng vào văn phòng.

``Có chuyện gì mà ồn ào thế kia--''

*RẦM!*

Cánh cửa bật tung ra, và hai bóng người lao thẳng vào phòng với tốc độ kinh người - đó là Korone và Towa.

``Đứng lại đó cho chị nào\~{}! Chúng ta chưa nói chuyện xong đâu!!'' Korone hét lên, giọng nói vẫn trông vui vẻ nhưng ẩn chứa đầy sự đe dọa, đôi mắt sáng rực lên.

``\textbf{Tha cho em đi, senpaiiiiii!!}'' Towa khóc thét, chân chạy bán sống bán chết, hai tay vẫn bị trói chặt như cũ.

Chỉ trong nháy mắt, hỗn loạn bùng nổ. Nàng Ác ma với đôi tay vẫn bị trói chạy thục mạng vào trong phòng, còn nàng Khuyển thì đuổi theo sát gót, đôi mắt sáng rực như đang săn mồi, chiếc lưỡi thè ra ngoài vì quá thích thú.

Ryuuhou vội vàng đứng dậy, lùi sang một bên để tránh, mắt mở to nhìn cảnh tượng trước mắt: ``Trời ơi, Korone-san vẫn chưa chịu thả Towa-san ra sao!?''

Chưa kịp ai đó hiểu chuyện gì đang xảy ra, một giọng nói bất ngờ vang lên từ cái hố dưới sàn nhà mà Pekora đã đào từ trước: ``Pe-ko-raaaaa! Mau xin lỗi tôi mau!!''

Mặt đất hơi rung lên, rồi từ từ, một mái tóc xanh lục cùng khuôn mặt đầy tức giận lộ ra. Trong tay Rushia là một con dao găm nhỏ (chắc chắn chỉ để dọa thôi chứ không có ý làm gì thật, hy vọng là vậy). Ánh mắt cô sắc lạnh, tràn đầy sự đe dọa.

\img{9-3l}

Lamy hoảng sợ tột độ, lập tức kéo ghế lùi thật xa, mặt tái nhợt giống hệt màu tóc của mình: ``C-C-Chờ đã! Mọi người có chuyện gì vậy!?'' cô kêu lên, giọng đầy hoảng loạn, tay vịn chặt vào mép bàn.

Cùng lúc đó, Korone vồ lấy Towa, cả hai cùng té ngã trên sàn nhà, tạo thành một trận giằng co ồn ào. Nàng Ác ma cố gắng vùng vẫy để thoát thân, lắc lư như một con cá mắc cạn, còn nàng Khuyển thì cười khúc khích, trông chẳng có vẻ gì là định buông tha đàn em, như thể đây là trò chơi yêu thích nhất của cô.

Rushia lúc này đã hoàn toàn trồi lên khỏi cái hố, chỉ thẳng con dao về phía Pekora, mặt đỏ bừng vì tức giận: ``Giải thích đi, con Thỏ kia!!''

``Đ-Đợi đã!! Tôi đâu có làm gì đâu peko!!'' Pekora vội vàng xua tay, cố gắng tự vệ, lùi dần về sau mỗi khi Rushia tiến thêm một bước.

Ryuuhou bước tới, đứng giữa Rushia và Pekora, giơ tay ra hiệu bình tĩnh: ``Rushia-san, bình tĩnh nào! Chị đang làm mọi người sợ hãi đấy! Em biết chị đang giận, nhưng dùng dao để giải quyết vấn đề không phải là cách đúng đâu!''

Cả căn phòng lập tức trở nên hỗn loạn, tiếng nói của người này xen lẫn tiếng la hét của người kia, tiếng thét của cô này hòa vào tiếng kêu của cô kia. Không khí trong phòng như một cơn lốc xoáy đang cuốn tất cả mọi người vào vòng xoáy hỗn loạn.

``M-Mọi người bình tĩnh lại đi mà!!'' Lamy mặt xanh tái, cố lao tới để ngăn cản mọi người, nhưng chỉ nhận lại những cú đẩy tay và những tiếng la hét.

Trong khi đó, Pekora đang cố gắng nép mình vào bàn tư vấn, tránh né những nhát dao mà Rushia đang tấn công dồn dập: ``Họp mặt bạo lực à peko!?'' cô hét lên, giọng đầy hoảng sợ.

Nhưng Rushia chẳng hề quan tâm. Cô chĩa dao về phía nàng Thỏ, gầm lên, mắt đỏ ngầu: ``\textbf{Còn không mau giải thích cho tôi nhanh lên!!}''

Và đúng lúc đó\dots

``{\Large \textbf{Thôi!!!}}'' Một tiếng quát mạnh mẽ vang vọng khắp căn phòng, khiến cả không gian dường như rung chuyển. Mọi tiếng động lập tức im bặt, ngay cả những va chạm trên sàn cũng như ngừng lại.

Korone và Towa ngừng vật lộn. Rushia đứng yên tại chỗ, con dao vẫn còn giơ lên, không thể nhúc nhích. Lamy và Pekora ngồi im thin thít. Ryuuhou cũng đứng lặng, tay vẫn giơ lên như một bức tượng.

Kuon thở dài, rồi nhìn sang bàn tư vấn của Lamy và Pekora, giọng nói vẫn bình thản như chưa từng có chuyện gì xảy ra: ``Hai người cứ tiếp tục nói chuyện đi.''

Ryuuhou từ từ hạ tay xuống, nhìn Kuon với vẻ đầy ngưỡng mộ: ``Onii lúc nào cũng giữ được bình tĩnh trong mọi tình huống thật nhỉ. Em học được bài học quý giá rồi đấy.''

Hai cô nàng nhìn nhau, rồi lại nhìn cậu, có vẻ vẫn chưa thể hiểu hết tình huống bất ngờ vừa xảy ra. Nhưng cuối cùng, họ cũng quay lại với việc chia sẻ những tâm tư của mình, như thể chẳng có gì xảy ra.

Ba người còn lại thì tròn mắt nhìn cậu, chẳng hiểu chuyện gì đang diễn ra.

Towa, với hai tay vẫn bị trói chặt, lên tiếng bất bình, giọng đầy phẫn nộ: ``Anh làm như thể bọn em bị bỏ rơi vậy!''

Rushia cũng không kém cạnh, vẫn cầm dao chỉ về phía Kuon: ``Anh không thấy chuyện này thật sự nghiêm trọng sao, Kuon-senpai!?''

Ngược lại với hai người, Korone tỏ vẻ vô tư như thể chuyện này chẳng liên quan đến mình, cười toe toét: ``Được đó, được đó\~{}'' cô nói, giọng như một đứa trẻ vừa kết thúc trò chơi đuổi bắt.

Ryuuhou nhìn ba người, lắc đầu: ``Các chị ơi, em nghĩ mọi người nên ngồi xuống và nói chuyện bình thường thôi. Cứ thế này thì đến tối cũng chẳng giải quyết được gì đâu.''

Kuon cũng lắc đầu, ngồi xuống ghế rồi nhìn ba người vừa gây ra hỗn loạn: ``Rồi, có chuyện gì mà ồn ào như chợ vỡ thế?''

Lập tức, cả ba người cùng nhao nhao lên, mỗi người nói một câu, chẳng ai nghe được lời của ai. Giọng nói chồng chéo lên nhau, tạo thành một bản hòa tấu đầy lộn xộn.

``Tên thỏ này--!!''

``Em đâu có làm gì sai đâu!!''

``Vật lộn với nhau vui lắm, anh có muốn thử không\~{}?''

Cứ thế, ba cô gái tiếp tục nói chen ngang vào nhau, giọng người này lấn át giọng người kia.

Ryuuhou khoanh tay, nhìn họ với vẻ bất lực nhưng không giấu được nụ cười: ``Các chị mà cứ thế này thì khác nào một buổi tập nhạc không có nhạc trưởng. Ai cũng muốn chơi phần của mình, nhưng chẳng ai chịu lắng nghe người khác cả.''

Kuon nhíu mày, giơ tay ra hiệu dừng lại, cả ba cô gái lập tức im lặng. Cậu nhìn từng người một, giọng nói đều đặn nhưng đầy uy lực: ``Từng người một nói với anh. Bắt đầu với Ru-chan.''

Rushia vẫn hậm hực, nhưng cũng hạ con dao xuống, bắt đầu kể lể mọi chuyện.

Vậy là, từ một buổi tư vấn đơn giản giữa Lamy và Pekora, giờ đây đã trở thành một buổi tư vấn mở rộng với tất cả mọi người tham gia. Một bên vẫn là hai người tiếp tục chia sẻ tâm tư, còn bên kia, Kuon với trợ lý Ryuuhou vô tình trở thành một ``bác sĩ tâm lý'' bất đắc dĩ, lắng nghe và giải quyết từng vấn đề lần lượt.

Ryuuhou ngồi xuống cạnh Kuon, thầm thì nhỏ: ``Onii, em thấy hôm nay anh làm quá nhiều việc rồi đấy. Vừa quản lý công ty, vừa làm tư vấn viên, lại còn kiêm cả cảnh sát nữa. Em nghĩ nếu anh mở trung tâm tư vấn thì chắc chắn sẽ ối khách hơn Pekora-san rất nhiều.''

Kuon chỉ thở dài, lắc đầu: ``Anh chỉ muốn về nhà thôi.''

``Chịu khó một chút đi Onii, sắp xong rồi,'' Ryuuhou cười, vỗ nhẹ vào vai anh.

\dots Hôm nay đúng là một ngày thật dài.

\ct

Rushia khoanh tay với vẻ vẫn còn hậm hực, hít vào một hơi thật sâu rồi mới bắt đầu trút bầu tâm sự. ``\textbf{Pekora nói em lép á!!}''

Kuon hơi nghiêng đầu, hỏi với giọng bình thản: ``Tại sao em lại nghĩ vậy?''

``Cô ấy bảo rằng `mỗi người sinh ra vốn đã là bản thân mình từ đầu', thế thì chẳng phải cô ấy đang ám chỉ em lép sao!?'' giọng Rushia pha lẫn sự ấm ức và bực bội, nhưng ẩn sâu trong đó còn có nỗi chua xót, như thể cô đã giữ kín nỗi tủi thân này trong một thời gian quá dài.

Ryuuhou ngồi cạnh anh trai mình, lắng nghe một cách thầm lặng với ánh mắt tràn đầy sự đồng cảm. Cô bé khẽ nghiêng đầu, giọng nói nhẹ nhàng nhưng vẫn pha chút hài hước: ``Rushia-san, em hiểu cảm giác mà chị đang trải qua. Có những lúc em cũng bị người khác trêu chọc vì ngoại hình của mình, dù em chẳng để lòng lắm.''

Cậu Tổng chống cằm, suy nghĩ lát rồi đáp lời một cách nhẹ nhàng: ``Anh cũng không nghĩ em ấy có ý gì xấu. Chỉ là em ấy không khéo ăn nói thôi.''

Dù đã có lời giải thích, nàng Pháp sư vẫn chẳng thấy thỏa mái. Thậm chí cô còn liếc nhìn Pekora - người đang ngồi co ro ở góc bàn - với ánh mắt như muốn bắn chết đối phương.

Cậu thở dài, đột nhiên hỏi với giọng nhẹ nhàng để không khiến cô nàng cảm thấy khó xử: ``Nói chuyện này hơi nhạy cảm một chút, nhưng trong nhà em có ai có vóc dáng `boing-boing' không? Em có thể giận anh sau cũng được.''

Rushia đột nhiên đứng im, ánh mắt dao động không ngừng. Cô mím chặt môi, trông vô cùng lưỡng lự khi phải đối mặt với câu hỏi khó trả lời này. Một lúc lâu sau, cô mới bối rối đáp nhỏ, má cô đỏ bừng lên: ``\hl{\dots Mẹ em có ạ-nanodesu.}''

Ryuuhou bật cười nhẹ, giọng nói đầy trêu chọc nhưng hoàn toàn không mang ý mỉa mai: ``Vậy thì Rushia-san còn hy vọng lắm đấy. Em có quen một người từng bảo rằng gen di truyền là một yếu tố rất quan trọng. Chỉ cần chị kiên nhẫn thôi, mọi thứ sẽ đến với chị.''

``Vậy em đã bao giờ hỏi mẹ mình làm thế nào để có được vóc dáng như bà ấy chưa?'' Kuon tiếp tục câu chuyện với giọng điềm tĩnh.

Rushia cúi đầu, ngón tay vô thức vẽ những đường lung tung trên đùi quần, như thể đang cố tìm một điểm tựa. Giọng nói của cô nhỏ dần vì sự ngại ngùng, lắp bắp: ``\hl{Em\dots{} có hỏi. Nhưng mẹ em chỉ bảo cứ từ từ, rồi ngực em sẽ lớn dần lên thôi.}''

Kuon bật cười nhẹ, lắc đầu: ``Đúng đó! Em sẽ có được vóc dáng boing-boing đó, chỉ cần chút thời gian thôi.''

Ryuuhou thêm vào với giọng như một người đã từng trải: ``Rushia-san, thời gian là thứ mà không ai có thể đẩy nhanh được. Em cũng đã học được điều đó khi\dots\ ờm, khi em thấy mọi người xung quanh mình thay đổi từng ngày một cách từ từ. Chị cứ để mọi thứ diễn ra tự nhiên là được.''

``Nhưng như vậy thì lâu lắm! Chị muốn có nhanh cơ!!'' Rushia tỏ ra bất mãn, giương đôi mắt đỏ long lanh về phía cậu với vẻ tuyệt vọng, y hệt một đứa trẻ đang đòi quà mà không chịu chờ đợi.

Kuon lặng lẽ quan sát cô một hồi, rồi đột nhiên đặt ra một câu hỏi khiến cô nàng đứng hình hoàn toàn: ``Sao ngoại hình lại quan trọng đến thế với em?''

Rushia mở miệng định phản bác lại, nhưng không hiểu sao, lời nói của cậu lại như một tảng đá nặng đè lên tâm trí cô. Cô lập tức ngậm chặt miệng, đôi mắt đỏ hoe mở to nhìn hai anh em.

Quan trọng ư? Chẳng lẽ nó lại không quan trọng sao? Nhưng\dots

``Tại sao em lại muốn mình có vóc dáng phổng phao hơn? Anh hỏi thật đấy,'' cậu vẫn giữ vẻ điềm nhiên, kiên nhẫn chờ câu trả lời từ cô nàng, giọng nói chẳng hề mang sắc thái phán xét.

Rushia bối rối, ánh mắt dao động liên tục. Cuối cùng, cô thở hắt ra, thú nhận nhỏ giọng như thể đang thú nhận một tội lỗi: ``\dots Tại fan của em chê em giống cái thớt\dots'' giọng cô nghẹn lại ở cuối câu, nỗi đau ấy dường như vẫn còn rất mới mẻ.

Kuon khẽ gật đầu, đưa ra kết luận thẳng thắn nhưng vẫn đầy sự thấu hiểu: ``Tức là em đang ghen tị với những thành viên khác trong nhóm, phải không?''

Đến đây, nàng Pháp sư đột nhiên cứng họng. Cô chẳng thể nào phản bác lại lời cậu nói. Câu hỏi đó như trúng tim đen, chạm đúng vào cái tôi và lòng tự trọng của chính cô. Cô cắn chặt môi, chẳng biết phải trả lời thế nào.

Ryuuhou nhẹ nhàng đặt tay lên bàn, giọng nói tràn đầy sự đồng cảm: ``Rushia-san, em không rõ cảm giác bị chê bai là như thế nào, nhưng em biết rằng luôn có những người yêu quý chị vì con người thật của chị. Họ ở bên cạnh chị không phải vì ngoại hình, mà vì sự ấm áp mà chị mang lại cho họ.''

``Không có ý xúc phạm đâu, nhưng anh hiểu cái cảm giác đó,'' Kuon tiếp lời. ``Thường thì mọi người sẽ lấy oshi của mình ra so sánh với người khác, rồi đem khuyết điểm ra để trêu chọc. Anh cũng vậy, có vô số lỗi lầm.''

Rushia ngước lên, trông có chút ngạc nhiên, như thể không ngờ người đàn ông trước mặt lại có thể hiểu cô đến vậy: ``\dots Thật vậy ạ?''

``Ừ. Em cứ nhìn mấy đứa hay stream cùng anh đi. Lúc nào họ chả trêu là: `Sao cậu nói nhanh thế', rồi `Thằng này chơi game dở vãi, còn đếch bằng oshi nhà mình', còn rất nhiều câu nói vùi dập tương tự. Anh bị nhiều rồi nên hiểu rõ. Vấn đề không phải là họ nói gì, mà là cách mình đối mặt với những lời đó.''

Ryuuhou cũng gật gù đồng tình, bổ sung thêm: ``Em cũng từng bị nhận xét là quá thẳng thắn, có lúc nói chuyện mà khiến người khác hiểu sai ý. Nhưng em đã học được cách không để những lời đó ảnh hưởng đến mình quá nhiều. Rushia-san, chị cũng có thể làm được điều đó.''

Rushia im lặng hoàn toàn.

Cô nhớ lại những lần bị người hâm mộ gọi là ``thớt,'' những lúc cô bị đem kích thước vòng một ra so sánh với Marine hay Noel. Cô ghét cái cảm giác đó. Cô cảm thấy mình bị xem nhẹ, bị đánh giá thấp hơn người khác. Nhưng\dots\ liệu có phải tất cả mọi người đều chê bai cô theo hướng tiêu cực không?

Kuon dừng lại một chút, rồi tiếp tục nói với giọng trầm ấm, như một người anh đang động viên em gái của mình: ``Đôi khi người ta trêu em là `thớt' cũng chỉ vì họ quý em thôi. Nếu không quý thì họ còn chả thèm xem stream của em, nói gì đến việc ngồi lại bình luận.''

Ryuuhou mỉm cười, nói thêm: ``Rushia-san, em có một người quen, cũng hay bị trêu về ngoại hình. Nhưng cô ấy từng nói rằng những lời trêu đó đều xuất phát từ sự thân thiết, và nếu không có những người đó thì cuộc sống của cô ấy sẽ buồn đi rất nhiều. Chị thử nghĩ xem, những người hay trêu chị nhất có phải cũng là những người thân thiết nhất với chị không?''

Lời nói này khiến Rushia mở to mắt, nhưng cô vẫn chẳng thể nào thốt lên được một câu nào. Cô tiếp tục lắng nghe những lời tư vấn từ Kuon và Ryuuhou. ``Ai cũng có khuyết điểm của mình, nhưng cách mà chúng ta giải quyết nó mới là điều quan trọng nhất,'' cậu nói với giọng đầy chắc chắn.

Cậu Tổng khẽ nhếch môi, rồi đột nhiên cất giọng mạnh mẽ, như một lời tuyên bố đầy quyền lực: ``\textbf{Anh chả cần biết các em `thớt' hay là `bưởi' hay là củ quặc gì, chỉ cần mấy đứa làm việc tốt, đóng góp được cho công ty là ổn rồi!}''

Lời nói chắc nịch ấy, giọng cậu nói mạnh đến mức khiến cả căn phòng dường như rung chuyển nhẹ. Dù không hề có ý làm ai sợ hãi, nhưng áp lực từ trong lời nói của cậu lại rõ ràng hơn bao giờ hết. ``\textbf{Cái quan trọng là cái `gỗ' bên trong, chứ không phải lớp `nước sơn' ở bên ngoài! Chỉ có thế thôi!!}''

Nàng Pháp sư tròn mắt nhìn cậu Tổng, cảm thấy như có một luồng điện chạy xuyên qua người.

Không hiểu sao\dots\ cô lại cảm thấy tim mình đập nhanh hơn một chút.

Ryuuhou nhìn Rushia, rồi nói với giọng nhẹ nhàng nhưng ấm áp: ``Rushia-san, em thấy những gì Onii nói đều đúng hết. Nếu chị có thể tự tin hơn vào bản thân, thì những lời chê bai kia sẽ chẳng thể nào làm tổn thương được chị. Và em tin rằng chị đã có đủ tình yêu thương từ những người xung quanh để làm được điều đó.''

Trong mấy năm qua, Rushia luôn mang trong mình sự tự ti về ngoại hình. Cô từng nghĩ rằng nếu mình không có điểm gì nổi bật, nếu không đủ quyến rũ như những thành viên khác trong nhóm, thì một ngày nào đó những người hâm mộ sẽ rời bỏ cô. Cô luôn sợ hãi điều đó, ngày ngày lo lắng viễn cảnh ấy sẽ xảy ra.

Nhưng\dots

Nếu ngoại hình thực sự quan trọng đến vậy, thì làm sao cô có thể đạt được vị trí như hiện tại? Làm sao cô có thể có được những người hâm mộ trung thành, những người luôn ủng hộ cô hết mình?

Lời nói của hai anh em nhà Hikawa mang trọng lượng hơn bất cứ điều gì cô từng nghe trước đây.

Cuối cùng, cô hạ ánh mắt xuống, ngồi im lặng suy ngẫm lại tất cả. Đôi tay cô vô thức siết chặt lại, như thể đang nắm giữ một quyết định vô cùng quan trọng.

Thấy vậy, Kuon cũng dịu giọng lại, không muốn tạo thêm áp lực cho cô: ``\dots Xin lỗi em nếu anh vừa rồi hơi gắt. Cứ suy nghĩ thoải mái đi. Anh cho em thời gian.''

Ryuuhou nhìn Rushia với ánh mắt đầy sự động viên: ``Rushia-san, chị cứ suy nghĩ kỹ nhé. Em tin rằng chị sẽ tìm ra câu trả lời phù hợp nhất với bản thân mình.''

Nàng Pháp sư vẫn ngồi yên đó, chìm trong suy nghĩ, đôi mắt nhìn vào khoảng không vô định như thể vừa nhìn thấy một chân trời mới mở ra trước mắt.

Kuon quay sang nhìn hai người còn lại - Towa và Korone - rồi nói với giọng đã trở lại bình thường: ``Được rồi, người tiếp theo!''

Ryuuhou cũng nhìn sang Towa, mỉm cười: ``Towa-san, đến lượt chị rồi. Có vấn đề gì mà bọn em có thể giúp đỡ không?''

\ct

Đến lượt mình, Towa im lặng bước đến ghế và ngồi xuống, ánh mắt lướt qua một chút do dự, những vết tích của sợi dây thừng vẫn còn lưu lại trên cổ tay dù đã được tháo ra từ lúc nào. Trước khi cuộc trò chuyện chính thức bắt đầu, Kuon nhẹ nhàng giúp cởi bỏ hoàn toàn những nút thắt còn vướng lại trên tay cô. Nàng Ác ma thở phào nhẹ nhõm, duỗi thẳng các ngón tay để xoa dịu cảm giác tê mỏi, rồi lén liếc sang Korone, người vẫn đang ngồi nhìn cảnh tượng với vẻ thích thú, như thể vừa chơi xong một trò chơi cực kỳ thú vị. Towa khẽ hừ một tiếng trong miệng, ánh mắt lướt qua sự ngán ngẩm và trách móc, như đang nhắn nhủ: ``Lần sau mà còn dám dọa em như thế nữa thì coi chừng, Korone-senpai!''

Kuon hắng giọng một tiếng, kéo sự chú ý của cô quay lại với cuộc nói chuyện. Cô nàng giật mình, nhíu mày nhìn cậu Tổng. Dù đã chuẩn bị tinh thần từ trước, nhưng khi đối diện trực diện với ánh mắt nghiêm nghị của cậu Tổng, cô vẫn không tránh khỏi cảm giác hơi căng thẳng. Bên cạnh, Ryuuhou khẽ chống cằm, ánh mắt nhìn Towa tràn đầy tò mò nhưng vẫn giữ được nụ cười dịu dàng trên môi.

Sau một lúc do dự chần chừ, nàng Ác ma mới dần mở lòng trình bày vấn đề của mình. Giọng nói cô nhỏ nhẹ, vẫn còn vương chút ngập ngừng, vì đây không phải là lần đầu tiên cô chia sẻ mong muốn này với Kuon, kể từ lần đầu tiên cô đặt chân đến văn phòng công ty: ``\dots Em muốn mình trở nên giống một ác ma hơn ạ.''

Kuon chỉ khẽ gật đầu, vẻ mặt không hề có chút ngạc nhiên nào. Ryuuhou cũng khẽ nghiêng đầu, như đang suy ngẫm về điều gì đó, nhưng vẫn giữ im lặng để anh trai mình tiếp tục cuộc nói chuyện. ``Cái này anh biết. Ngay từ ngày đầu tiên em mới vào công ty, em cũng đã từng nói với anh điều tương tự. Anh còn nhớ lúc đó anh đã bảo em cứ mạnh dạn làm những trò chơi khăm để thỏa mãn niềm đam mê, đúng chứ?''

``Dạ\dots vâng ạ.''

``Anh cũng biết em thích chơi khăm mà. Nhưng có điều anh vẫn thắc mắc mãi, nếu là một ác ma đích thực như em thì tại sao lại có phần nhát gan thế kia?''

Câu hỏi bất ngờ khiến Towa khựng lại hoàn toàn. Đôi mắt xanh lá mạ trong trẻo của cô thoáng chốc dao động khó hiểu. Dù biết đây không phải là lần đầu Kuon đặt ra câu hỏi này, cô vẫn rơi vào trạng thái lúng túng, chẳng biết phải mở lời giải thích thế nào. ``\dots C-Cái đó thì\dots''

Ryuuhou nhẹ nhàng lên tiếng, giọng nói vừa đủ lớn để không làm gián đoạn cuộc nói chuyện, mà còn giúp Towa cảm thấy bớt cô đơn hơn: ``Towa-san, em nghĩ việc em muốn trở nên ác ma hơn nhưng lại luôn sợ sẽ làm tổn thương đến người khác chính là điều làm nên sự đặc biệt của chị. Không phải ai cũng có thể giữ được sự cân bằng tuyệt vời đó đâu.''

Chưa để Towa kịp phản hồi gì thêm, Kuon lại tiếp tục: ``Em có một trái tim rất nhân hậu, đúng chứ? Và tất cả những trò chơi khăm mà em tạo ra cũng chỉ để mang lại tiếng cười cho mọi người, chứ chẳng bao giờ có ý định làm hại ai, đúng không?''

``\dots!'' cô nàng mở to đôi mắt, nhìn chằm chằm vào Kuon như thể vừa bị bóc trần hoàn toàn những suy nghĩ thầm kín nhất trong lòng.

Cậu Tổng khẽ bật cười, rồi lại hỏi thêm một câu hỏi sâu sắc: ``Đã bao giờ em tự hỏi bản thân rằng `Thật sự mình là ai?' chưa?''

``\dots Em chưa bao giờ nghĩ đến chuyện đó,'' nàng Ác ma đáp nhỏ giọng, lí nhí như đang nói với chính mình.

Nàng Trưởng nhóm cũng chia sẻ thêm, giọng nói như đang kể lại một kinh nghiệm cá nhân đã trải qua: ``Em cũng từng có những lúc tự hỏi mình là ai và mình thực sự muốn gì. Nhưng rồi em dần nhận ra rằng đôi khi câu trả lời không nằm ở việc phải chọn một trong hai mặt đối lập, mà là chấp nhận và dung hòa cả hai phần đó.''

``Nếu em tìm ra được câu trả lời cho chính mình, em sẽ biết phải làm gì tiếp theo,'' cậu Tổng tiếp lời em gái mình. ``Còn một điều nữa: em có bao giờ tự hỏi mình `Mình làm tất cả những điều này vì bản thân, hay là vì một lý do nào khác không?' thế nào chưa?''

Lần này, Towa không trả lời ngay lập tức. Đôi bàn tay cô khẽ siết chặt vào đầu gối, ánh mắt cụp xuống nhìn sàn nhà, như thể đang miệt mài đào sâu vào những suy nghĩ mà từ trước đến nay cô chưa bao giờ dám đối mặt.

``Vì gia đình em mong muốn em làm như vậy, phải không?'' Kuon hỏi, khẽ nhướng một bên mày.

Towa giật mình đến mức cái đuôi nhỏ phía sau lưng cô dựng đứng thẳng lên, như vừa bị ai đó tấn công một cách trực diện và bất ngờ. ``S-Sao anh\dots\ Sao anh lại biết được chuyện đó!?'' giọng cô tràn đầy sự kinh ngạc, như thể một bí mật lớn nhất của mình vừa bị vạch trần trước mặt mọi người.

Ryuuhou khẽ mỉm cười, giọng nói nhẹ nhàng an ủi: ``Onii của em có khả năng nhìn người rất tốt đấy, Towa-san. Em cũng đã học được rất nhiều điều từ cách quan sát của anh ấy.''

``Anh có nhiều người bạn từng rơi vào hoàn cảnh tương tự như em, nên anh hiểu rõ cảm giác đó,'' Kuon nhẹ nhàng giải thích. ``Áp lực từ gia đình, sự kỳ vọng của bố mẹ đặt lên con cái để con có được những điều tốt đẹp nhất, anh đều hiểu cả. Bản thân anh cũng đã từng trải qua những điều tương tự như em. Nhất là khi em gái nhà anh đã bắt đầu kiếm được tiền trong khi anh vẫn còn ăn bám bố mẹ.''

Towa im lặng trong một lúc khá dài, rồi mới nhẹ nhàng gật đầu đồng tình. Cô thở ra một hơi thật sâu, như thể đang chuẩn bị tinh thần để mở lòng về một câu chuyện chưa bao giờ cô kể với bất kỳ ai: ``Không giấu gì anh, từ ngày đầu tiên bước chân xuống Trái Đất, em vẫn luôn giữ được một chút bản chất ác ma trong người. Gia đình ở nhà gửi tiền sinh hoạt cho em để em có thể học hỏi và thích nghi với thế giới mới này, nhưng\dots\ kể từ khi em biết đến trò chơi điện tử, em dần dần sa đà vào đó không thể dừng lại. Rồi cứ thế mà quên đi\dots''

Ryuuhou vẫn giữ nụ cười ấm áp, như thể đã phần nào hiểu được cảm giác của Towa: ``Towa-san, em nghĩ việc chị tìm thấy niềm vui đích thực trong game cũng là một cách để chị khám phá thế giới này. Có thể đó không phải là con đường mà gia đình mong đợi từ đầu, nhưng nó lại chính là con đường mà chị đã chọn cho riêng mình.''

``\dots Và rồi, em đã quên mất hoàn toàn mục đích ban đầu khi mình đến đây, đúng không?'' Kuon khẽ hỏi.

``\dots''

Towa khẽ cười, nhưng trong đáy mắt lại chứa đầy sự bối rối và mâu thuẫn. ``Vâng\dots\ Nhưng em cũng dần nhận ra rằng thế giới loài người không hề giống như những gì gia đình em từng dạy dỗ. Nó hoàn toàn trái ngược, hoàn toàn khác biệt.''

Kuon bật cười nhẹ nhàng trước lời bộc bạch dù hơi vụng về nhưng vô cùng chân thành từ cô nàng. ``Vậy theo em, con người thực sự ở đây là như thế nào?''

Towa suy nghĩ một lúc, rồi mới chậm rãi trả lời, giọng nói dần dần trở nên ấm áp hơn: ``Con người ở đây đều rất tốt bụng. Họ khác xa hoàn toàn so với những gì em từng nghe kể về loài người. Họ còn rất đoàn kết, lúc nào cũng sẵn sàng giúp đỡ lẫn nhau. Và ở đây có rất nhiều thứ thú vị, những trò giải trí mà ở quê nhà em không bao giờ có được.''

Cô ngừng lại một chút, rồi tiếp tục với giọng nói càng lúc càng nhỏ đi, như thể đang thú nhận một điều gì đó cực kỳ riêng tư: ``Quan trọng nhất là\dots\ họ đều rất yêu quý em.''

Kuon khẽ gật đầu ghi nhận. ``Ừ, tiếp đi em.''

``Vì vậy, để đáp lại lòng tốt của mọi người, em đã quyết định gia nhập \hll. Chính nhờ sự chỉ dẫn và giúp đỡ của mọi người, em mới có thể ngồi ở đây được như ngày hôm nay ạ.''

Kuon vẫn giữ im lặng lắng nghe, không bao giờ ngắt lời Towa khi cô đang mở lòng. Ryuuhou cũng vậy, chỉ khẽ mỉm cười với Towa như một lời động viên thầm lặng.

``Trong mỗi buổi stream, em luôn cố gắng hết sức để pha trò, trêu chọc mọi người nhằm mang lại tiếng cười cho khán giả. Và rồi\dots''

Towa bĩu môi đầy sự uất ức, đến mức bản thân cô không hay biết giọng nói đã lớn dần lên lúc nào: ``Những người hâm mộ lại gọi em là thiên thần. \textbf{THIÊN THẦN!} Trong khi em đâu có phải là Kanata đâu! Anh có cảm thấy chuyện này thật khó chịu không!?''

Ryuuhou bật cười nhẹ, nhưng hoàn toàn không phải với ý định trêu chọc Towa: ``Towa-san, em nghĩ fan của chị gọi chị là thiên thần đơn giản là vì họ thấy được sự tốt bụng và ấm áp trong con người chị. Đó đâu phải là điều xấu đâu chứ.''

Câu than thở bất lực của Towa đã phần nào làm bầu không khí căng thẳng trong phòng vơi đi. Kuon cũng khẽ mỉm cười, cậu đã dần hiểu rõ hơn về nàng Ác ma đang lạc lối trước mặt mình.

Kuon khoanh tay trước ngực, ánh mắt sắc bén nhưng không hề có ý trách móc ai. Cậu khẽ gật đầu, như thể cuối cùng cũng đã thấu hiểu được tất cả những gì Towa mong muốn: ``Vậy là anh đã hiểu tại sao em lại muốn mình trở nên ác ma hơn rồi.''

``Dạ?''

Towa ngơ ngác nhìn cậu Tổng, đôi mắt chớp chớp liên tục, rõ ràng là không hiểu cậu đang ám chỉ điều gì. Cậu nhìn cô với vẻ mặt hoàn toàn nghiêm túc, không có một chút thay đổi nào, rồi mới nói một cách ngắn gọn: ``Tính cách của em\dots\ dường như đang đi ngược lại với bản chất mà em khao khát. Nếu em thực sự muốn trở thành một ác ma đích thực, thì em buộc phải từ bỏ đi cái trái tim tốt bụng mà em đang sở hữu.''

Những lời nói đó như một cú đánh mạnh vào tâm trí Towa. Cô không bao giờ ngờ rằng câu trả lời sẽ đến sớm và trực diện như vậy. Cảm giác bị bóc trần hoàn toàn khiến trái tim cô đập nhanh hơn bao giờ hết. Cô lắp bắp không nên lời, môi mấp mãi như thể không muốn chấp nhận sự thật đó. Towa trợn tròn mắt, cả người cô sững sờ lại như bị đóng băng tại chỗ.

``Nhưng mà\dots!'' cô bật dậy theo phản xạ tự nhiên, hai tay bấu chặt vào mép bàn. Chính phản xạ đó cũng chứng minh rằng đó không bao giờ là điều mà cô thực sự muốn làm.

Kuon vẫn giữ được vẻ điềm tĩnh vốn có, không hề nao núng trước phản ứng của nàng Ác ma, mà tiếp tục nói: ``Không thể làm được đúng không? Vì chính em cũng đâu có muốn phải từ bỏ bản thân mình như thế.''

Towa mở to đôi mắt, cảm giác mọi thứ bỗng chốc sáng tỏ ra như một cơn mưa rào mùa hạ vừa làm tan biến toàn bộ lớp sương mù dày đặc trong đầu. Cô vẫn giữ im lặng, chẳng biết phải nói thêm gì nữa. Cảm giác bất lực và bối rối dâng tràn trong lòng, nhưng cô lại không thể phủ nhận bất kỳ một lời nào mà Kuon vừa nói ra.

Ryuuhou khẽ nói thêm, giọng nói như một người đã từng trải qua cảm giác tương tự: ``Towa-san, em cũng từng có những lúc muốn biến mình thành một người hoàn toàn khác. Nhưng em đã học được rằng, đôi khi điều tốt nhất chúng ta có thể làm là chấp nhận bản thân và tìm cách dung hòa những phần khác biệt trong con người mình.''

``\dots''

``Bản chất của em đã là như vậy từ đầu rồi. Nhưng hãy nghe anh nói: em vừa khao khát được trở thành một ác ma, nhưng đồng thời cũng muốn giữ được mối quan hệ tốt đẹp với thế giới loài người, đúng chứ?''

Towa không nói thêm một lời nào, chỉ có thể nhẹ nhàng gật đầu đồng tình. Những lời nói của Kuon như một đòn đánh mạnh vào những điều mà cô đã cố giấu kín bấy lâu nay.

``Dạ\dots'' Cô trả lời một cách yếu ớt, như thể đang cố gắng tìm kiếm sự chắc chắn trong chính những suy nghĩ hỗn loạn của mình.

Kuon nhìn thẳng vào mắt cô, giọng nói vẫn điềm tĩnh nhưng chứa đựng nhiều sự sâu sắc: ``Vậy thì\dots\ Em chưa nhận ra rằng bản thân em đã trở thành một người như vậy rồi à?''

Towa vội quay mặt sang nhìn Kuon, đôi mắt tím của cô chợt nhòa đi vì xúc động. Cô không thể nào che giấu được sự bất ngờ lớn lao đó. Câu hỏi của Kuon như một tấm gương lớn soi rọi thẳng vào nội tâm sâu thẳm của cô. ``Ý-Ý anh là sao ạ, senpai?''

Kuon nhẹ nhàng giải thích, từng lời nói đều mang một sức nặng khó tả: ``Em có thể chơi khăm mọi người một cách vui vẻ mà không hề làm tổn thương ai, lại còn khiến tất cả mọi người đều yêu mến em. Em không phải là một ác quỷ thuần túy, cũng không hẳn là một thiên thần chính hiệu. Em là một cá thể vừa mang chút tinh nghịch của một ác ma, mà lại vẫn giữ được trái tim ấm áp của một con người.''

Ryuuhou cũng mỉm cười, khẽ gật đầu đồng tình với lời anh trai: ``Towa-san, đó chính là điều làm nên sự đặc biệt của chị. Chị chẳng cần phải buộc mình chọn một trong hai. Bởi lẽ chị đã là cả hai rồi.''

Towa mở to đôi mắt, nửa mừng nửa lo. Cô chưa bao giờ dám nghĩ mình có thể trở thành một con quỷ như vậy, nhưng những lời nói chân thành từ Kuon và Ryuuhou khiến cô cảm thấy bản thân mình thật sự có điều gì đó đặc biệt, một điều gì đó mà từ trước đến nay cô chưa từng nhận ra. Cô thậm chí còn không hay biết rằng mình chính là một trong những cá thể độc nhất vô nhị đó - một loại hình mà không phải ác ma hay bất kỳ ai cũng có thể sở hữu được.

``T-Thật vậy sao ạ\dots?'' đôi môi cô mấp máy khó nhọc, không thể thốt lên thêm bất kỳ lời nào. Cô khẽ siết chặt bàn tay vào ngực áo, như thể đang cố gắng cảm nhận rõ ràng nhịp tim đang đập mạnh trong lồng ngực của chính mình.

Kuon mỉm cười nhẹ nhàng, không hề vội vã, chỉ đưa ra một lời khuyên chân thành nhất. ``Ừ. Điều quan trọng nhất là phải luôn tin tưởng vào chính mình. Nhân tiện có Rushia-san ở đây nữa\dots''

Cậu chuyển ánh mắt sang Rushia, người từ nãy giờ vẫn ngồi im lặng, trầm tư theo dõi cuộc đối thoại với vẻ đăm chiêu. Kuon gọi cô nàng lại gần, đôi mắt vẫn giữ sự nghiêm chỉnh nhưng hoàn toàn không tạo ra áp lực nào. ``Anh cũng nhắn nhủ luôn với em: cứ là chính mình là tốt nhất. Tự làm chủ cuộc đời, tự định hướng con đường mình muốn đi. Đó chính là thứ làm nên giá trị của bản thân các em.''

Ngồi cạnh bên, Rushia không kìm được mà thở một hơi dài sau khi nghe lời cậu. Cô ngước nhìn Towa, mắt đã ngân ngấn lệ, giọng nói mang đầy niềm mong mỏi: ``Vậy\dots\ bọn em như thế này là đã ổn rồi hả anh?''

Kuon nhẹ nhàng gật đầu, nụ cười ấm áp của cậu làm tan đi bầu không khí căng thẳng còn sót lại: ``Đúng vậy. Các em hoàn toàn có thể hoàn thiện bản thân theo hướng tốt đẹp hơn, nhưng điều cốt lõi là phải làm những gì mà chính các em thực sự mong muốn.''

Hai cô gái nhìn nhau, nước mắt không còn kìm được nữa. Cùng lúc, cả hai lao tới ôm chầm lấy cậu, gục đầu trên vai cậu mà thổn thức. Towa siết chặt vòng tay như thể vừa tìm được người đầu tiên thực sự thấu hiểu mình, còn Rushia cũng không kém phần xúc động, nước mắt lướt dài trên má hồng.

``Cảm ơn anh\dots\ em chưa bao giờ nghĩ sẽ tìm được câu trả lời như thế này\dots!''

``Chưa từng có ai nói với em những điều chân thành như vậy cả\dots!''

Ryuuhou đứng một bên, mỉm cười ngắm nhìn khung cảnh ấm áp đó, giọng nói vừa hài hước vừa xúc động: ``Onii thật sự có năng khiếu làm tâm lý học đấy. Onii nên mở một văn phòng tư vấn cạnh Pekora-san luôn đi.''

Kuon thở dài, tay vô thức gãi gãi đầu ngọn, hơi lúng túng trước cái ôm đột ngột của hai cô gái. Cả hai vẫn chưa chịu nhả ra, khiến cậu càng thêm ngượng ngùng.

``Anh thấy mấy lời này cũng bình thường thôi. Người ở quê anh ai cũng nói như vậy cả,'' cậu đáp một cách khiêm tốn, chẳng hề nhận công sức của mình, mà chính sự khiêm tốn đó lại càng làm hai cô gái thêm cảm động.

Từ phía đối diện, Korone từ nãy giờ vẫn lặng lẽ quan sát, cái đuôi vẫy đều liên tục vì thích thú. Nhìn cảnh tượng xúc động kia, cô không nhịn được bật cười, tiếng cười trong trẻo như một chú cún con đang vui đùa.

Kuon liếc qua, nhíu mày hỏi: ``À, nhân tiện Koro-chan ở đây, có vấn đề gì muốn em tư vấn không?''

Korone nghiêng đầu, nụ cười tươi rói nở trên môi, cái đuôi vẫn vẫy nhịp nhàng: ``Không có gì đâu anh ạ! Em vừa chơi với Towa-chan xong, em vui lắm á!''

Cậu chớp mắt liên tục, toát cả mồ hôi hột: ``\dots Nghe câu đó hơi mờ ám nhỉ.'' Nhưng nàng Cún nâu chỉ cười khúc khích, chẳng có ý định giải thích gì thêm, cứ giữ bí mật nhỏ cho riêng mình.

Ryuuhou nhìn Korone, cười khẽ: ``Korone-san mà cứ thế này thì ai cũng sợ em mất. Nhưng em thích cái cách em chơi đùa với mọi người lắm.''

``Nhưng em không ngờ anh lại có thể làm nhà tư vấn đấy. Dường như anh cũng hiểu biết nhiều nhỉ?'' Korone nói, mắt vẫn híp lại vì niềm vui.

Kuon nhún vai, đáp gọn lỏn: ``Mấy chuyện kiểu này anh gặp mãi. Cũng bình thường thôi mà.''

Vậy là buổi tư vấn với ba người Pekora, Towa và Rushia đã thành công tốt đẹp. Vấn đề của mỗi người đều được gỡ nút theo một cách riêng. Bầu không khí căng thẳng ban đầu đã tan biến hoàn toàn, thay vào đó là niềm vui và sự ấm áp lan tỏa khắp phòng. Cả ba cô gái đều cảm thấy nhẹ nhõm hẳn, quay lại với công việc của mình với tinh thần phấn chấn hơn bao giờ hết.

Ryuuhou đứng dậy, vươn vai thỏa mái, rồi nhìn Kuon với ánh mắt đầy hài lòng: ``Onii, hôm nay anh đã làm rất tốt. Nếu anh có ý định mở trung tâm tư vấn thật sự, em xin làm trợ lý cho Onii luôn đấy.''

Kuon bật cười, vỗ vào lưng cô em gái: ``Cảm ơn, nhưng anh nghĩ có lẽ ta nên đóng cửa ngày hôm nay rồi.''

Mọi thứ cuối cùng cũng trở lại bình thường\dots

\dots cho đến khi\dots

\il

\pov{Hikawa Kuon}

\pov{Hikawa Kuon}

Sáng ngày hôm sau, tại ``Ký túc xá gia đình Hikawa'' - chính là căn nhà riêng của tôi - mọi thứ vẫn diễn ra trong nhịp điệu quen thuộc của một buổi sáng thường ngày.

Tôi vẫn giữ thói quen cũ, chuẩn bị đầy đủ đồ đạc để đến công ty. Bữa sáng vẫn vậy, ly cà phê vẫn vậy, chiếc áo khoác vẫn treo ở vị trí cũ. Mọi việc diễn ra đều đều như bao ngày khác, đến nỗi tôi chẳng hề hay biết rằng ngày hôm đó sẽ ẩn chứa một biến cố không ngờ tới.

Nhưng đúng lúc ấy, tiếng chuông điện thoại đột ngột vang lên, xé tan bầu không khí yên ả vốn có trong căn nhà. Nó vang lên giữa sự tĩnh lặng buổi sớm, giống như một tiếng còi báo động đánh thức mọi giác quan của tôi.

Tôi vươn tay cầm máy lên, nhấn nút nghe, giọng vẫn giữ bình thường dù trong lòng đã có chút bất an: ``Dạ, nhà Hikawa xin nghe ạ?''

Giọng nói của A-san vọng lại từ đầu dây, vẫn mang vẻ điềm nhiên như mọi lần, như thể chị ấy đang báo một tin tức thường nhật: ``Hikawa-san đúng phải không? Công ty ta vừa có chuyện lớn xảy ra rồi.''

Tôi đứng sững lại trong tích tắc. Lời nói tuy có vẻ thản nhiên, nhưng lại ẩn chứa một điều gì đó khiến lòng tôi dậy lên nỗi bất an khó tả. Chuyện lớn ư? Với \hll, chuyện lớn có thể là bất cứ điều gì, từ việc phát hành album mới cho đến vụ nổ văn phòng lần thứ bảy. Nhưng giọng A-san lần này có gì đó khác, một sự trầm lắng mà tôi không thể bỏ qua.

``Sao giọng chị vẫn bình thản thế này\dots?'' tôi nhíu mày hỏi lại, cảm thấy có chút khó chịu trong lòng, như thể đang đứng trước một cánh cửa mà tôi không chắc có muốn mở hay không.

A-san chỉ cười khúc khích, giọng vẫn giữ vẻ nhẹ nhàng: ``Tôi lúc nào chẳng vậy, cái cậu này! Nghe kỹ đây\dots''

Giọng nói của chị vẫn không thay đổi chút nào, thế nhưng từng lời thoát ra lại như những mũi kim nhọn đâm thẳng vào đầu tôi. Mỗi âm tiết đều rơi xuống với một sức nặng mà tôi không thể nào lường trước được. Tôi đứng hình hoàn toàn, bàn tay vô tình siết chặt lấy chiếc điện thoại đến mức các đốt ngón tay trở nên trắng bệch.

``C-Công ty của chúng ta\dots\ có người sắp phải tốt nghiệp rồi sao!?''

Cổ họng tôi như bị nghẹn lại, trái tim đập mạnh như bị bóp nghẹt, đầu óc trở nên trống rỗng trong mấy giây ngắn ngủi. Tôi không thể tin vào những gì mình vừa lắng nghe. Chúng ta đã cùng nhau vượt qua biết bao thăng trầm, cùng nhau đi qua những giai đoạn gian khó nhất. Tất cả vẫn còn đang dang dở\dots\ vậy mà sao lại đến lúc này?

Tôi chớp mắt liên tục, cố gắng trấn an bản thân, nhưng nỗi buồn về sự mất mát vẫn dâng lên từng chút một. Tôi tự hỏi trong lòng, liệu có phải điều này đã là điều đáng lẽ tôi phải đoán trước từ lâu rồi sao?

Ánh mắt tôi vô tình hướng lên bậc cầu thang dẫn tới tầng bảy, nơi căn phòng đặc biệt kia được hoàn thành vào cuối tháng Năm. Tôi còn nhớ rõ mồn một hôm đó, khi tôi, Ryuuhou, bố tôi và Aloe cùng nhau chung tay hoàn thiện nó – những tiếng cười, những giọt mồ hôi, và cả cảm giác hạnh phúc khi nhìn thấy thành quả sau bao ngày làm việc. Một căn phòng riêng biệt, nơi mà vị sếp lớn của tôi từng nhắc đến sẽ có người mới chuyển đến ở.

Lúc đó, tôi chẳng hề đặt nhiều câu hỏi về việc sắp xếp đó. Tôi chỉ đơn giản nghĩ đó là một sự điều động nhân sự thông thường, một bước đi trong kế hoạch phát triển của công ty. Nhưng bây giờ, sau khi biết được tin này\dots

Liệu đó có phải là những dấu hiệu báo trước về chuyện hôm nay không?

Tôi không dám chắc chắn về điều đó.

Nhưng có một điều tôi thấu hiểu rõ ràng: một trang mới trong cuộc hành trình của chúng ta sắp được mở ra. Và điều đó cũng đồng nghĩa, một trang cũ đang dần đến lúc phải khép lại.

\end{document}